\documentclass[12pt,a4paper]{article}

% Russian language support
\usepackage[T2A]{fontenc}
\usepackage[utf8]{inputenc}
\usepackage[russian]{babel}

% Math and figures
\usepackage{amsmath, amssymb}
\usepackage{graphicx}
\usepackage{float}

% Tables
\usepackage{tabularx}
\usepackage{booktabs}
\usepackage{array}
\usepackage{caption}
\captionsetup[table]{justification=raggedleft,singlelinecheck=false}
\setlength{\tabcolsep}{4pt}

% Hyperref and table of contents
\usepackage[hidelinks]{hyperref}
\usepackage{tocloft}

% Code listings
\usepackage{listings}
\usepackage{fancyvrb}
\lstset{
    basicstyle=\ttfamily\small,
    breaklines=true,
    frame=single,
    literate=
      {—}{{--}}1
      {←}{{$\leftarrow$}}1
      {→}{{$\rightarrow$}}1
}

% Geometry and spacing
\usepackage{geometry}
\geometry{left=2.5cm,right=2.5cm,top=2cm,bottom=2cm}
\setlength{\parindent}{1.25cm}
\setlength{\parskip}{0pt}

% For customizing section headings
\usepackage{titlesec}
\titleformat{\section}{\large\bfseries}{\thesection}{1em}{}
\titleformat{\subsection}{\normalsize\bfseries}{\thesubsection}{1em}{}

\begin{document}

% Title page
\begin{titlepage}
    \centering
    {\large \textbf{САНКТ-ПЕТЕРБУРГСКИЙ ГОСУДАРСТВЕННЫЙ УНИВЕРСИТЕТ}} \\[0.5cm]
    {\large Факультет прикладной математики--процессов управления} \\[0.5cm]
    {\large Программа бакалавриата} \\[0.2cm]
    {\large \textbf{``Большие данные и распределенная цифровая платформа''}} \\[4.5cm]
    
    {\Large \textbf{ОТЧЁТ}} \\[0.3cm]
    {\large По лабораторной работе № 1} \\[0.2cm]
    {\large По дисциплине ``Алгоритмы и структуры данных''} \\[0.2cm]
    {\large На тему ``Исследование задач глобальной оптимизации через генетический и роевой алгоритмы с модификациями.''} \\[4cm]
    
    \begin{flushright}
        Студент гр. 24Б15-пу \\
        Телицин А.Д. \\[0.5cm]
        Преподаватель \\
        Дик А.Г. \\[1cm]
    \end{flushright}
    
    \vfill
    \centering
    Санкт-Петербург \\
    2026 г.
\end{titlepage}

\tableofcontents
\newpage

\section{Цель работы}
\label{sec:goal}

Целью работы является сравнительный анализ эффективности генетического алгоритма с разными типами кодирования генотипа (вещественный и бинарный) и алгоритма роевого интеллекта в задачах глобальной оптимизации.

\section{Описание алгоритмов и постановка задачи}
\label{sec:algorithms}

\textbf{Генетический алгоритм} -- это эвристический алгоритм поиска, используемый для решения задач оптимизации и моделирования путём случайного подбора, комбинирования и вариации искомых параметров с использованием механизмов, аналогичных естественному отбору в природе.

Генетический алгоритм является разновидностью эволюционных вычислений, с помощью которых решаются оптимизационные задачи с использованием методов естественной эволюции, таких как наследование, мутации, отбор и кроссинговер.

Отличительной особенностью генетического алгоритма является акцент на использование оператора ``скрещивания'', который производит операцию рекомбинации решений-кандидатов, роль которых аналогична роли скрещивания в живой природе.

\textbf{Роевой алгоритм} (Particle Swarm Optimization -- PSO) -- это эвристический метод глобальной оптимизации, вдохновлённый социальным поведением живых организмов, таких как стаи птиц, косяки рыб или рои насекомых. Алгоритм моделирует движение группы агентов (частиц) в пространстве поиска, где каждая частица корректирует свою траекторию, ориентируясь на собственный лучший опыт и на опыт всей группы.

Роевой алгоритм является одним из подходов роевого интеллекта и широко применяется для решения многомерных задач оптимизации, в том числе там, где традиционные методы неэффективны.

Отличительной особенностью роевого алгоритма является использование коллективного обмена информацией: каждая частица хранит свою лучшую найденную позицию (личный оптимум) и учитывает глобальную лучшую позицию, обнаруженную всем роем. Благодаря этому обеспечивается баланс между исследованием (exploration) новых областей пространства и эксплуатацией (exploitation) уже найденных перспективных зон. Дополнительно вводится ограничение скорости частиц, предотвращающее чрезмерно большие перемещения и стабилизирующее процесс сходимости.

В работе исследуется работа данных алгоритмов в следующей задаче оптимизации:

\begin{equation}
f(x) = \frac{\sum_{i=1}^{n} \left( x_i^{4} - 16x_i^{2} + 5x_i \right)}{2} \rightarrow \min
\end{equation}

В работе будет рассмотрен случай \(n = 2\) в целях упрощения визуализации движения частиц в процессе работы алгоритма.

\section{Пошаговое описание работы ГА и РА и их блок-схемы.}

\subsection{Генетический алгоритм.}

\textbf{Пошаговое описание (версия с вещественным кодированием):}
\begin{enumerate}
    \item \textbf{Инициализация:} Создание начальной популяции из заданного числа особей. Каждая особь представляет собой вектор вещественных чисел, генерируемый случайным образом в заданном диапазоне.
    \item \textbf{Оценка приспособленности:} Вычисление значения целевой функции \(f(x)\) для каждой особи.
    \item \textbf{Запись лучшего решения:} Поиск особи с наименьшим значением функции в текущей популяции и сохранение её как лучшего решения, если оно лучше ранее найденного.
    \item \textbf{Формирование новой популяции.} Пока новая популяция не будет заполнена, выполняются шаги 4.1--4.4:
        \begin{enumerate}
            \item \textbf{Отбор (турнирный):} Случайным образом выбирается \(k\) особей из текущей популяции. Особь с наименьшим значением функции становится первым родителем. Процесс повторяется для выбора второго родителя.
            \item \textbf{Кроссовер (арифметический):} С заданной вероятностью применяется арифметический кроссовер: потомок вычисляется как взвешенная сумма родителей. Если кроссовер не применяется, потомок копируется от первого родителя.
            \item \textbf{Мутация (гауссова):} Для каждой координаты потомка с заданной вероятностью добавляется случайное возмущение из нормального распределения \(N(0, \sigma)\). Полученные координаты обрезаются до границ заданного диапазона.
            \item \textbf{Добавление потомка:} Полученная особь добавляется в новую популяцию.
        \end{enumerate}
    \item \textbf{Замена поколения:} Текущая популяция заменяется новой.
    \item \textbf{Проверка условия останова:} Если достигнуто максимальное число поколений, алгоритм завершается, иначе переход к шагу 2.
\end{enumerate}

\textbf{Пошаговое описание (версия с двоичным кодированием):}
\begin{enumerate}
    \item \textbf{Инициализация:} Создание начальной популяции из заданного числа особей. Каждая особь представляет собой закодированный через код Грея вектор вещественных чисел, генерируемый случайным образом в заданном диапазоне.
    \item \textbf{Оценка приспособленности:} Перевод закодированного кодом Грея числа обратно в вещественный вид. Вычисление значения целевой функции \(f(x)\) для каждой особи аналогично варианту с вещественным кодированием.
    \item \textbf{Запись лучшего решения:} Поиск особи с наименьшим значением функции в текущей популяции и сохранение её как лучшего решения, если оно лучше ранее найденного.
    \item \textbf{Формирование новой популяции.} Пока новая популяция не будет заполнена, выполняются шаги 4.1--4.4:
        \begin{enumerate}
            \item \textbf{Отбор (турнирный):} Для каждого родителя случайным образом выбирается заданное число индексов. Из всей популяции берется особь с наименьшим значением целевой функции, выбранная особь становится родителем.
            \item \textbf{Кроссовер:} С заданной вероятностью применяется одноточечный кроссовер для каждой координаты: выбирается точка разрыва, слева от которой вставляются гены первого родителя, справа -- второго. Если кроссовер не применяется, потомки являются копиями родителей.
            \item \textbf{Мутация:} Для каждого бита потомка с заданной вероятностью выполняется инверсия.
            \item \textbf{Добавление потомка:} Полученная особь добавляется в новую популяцию.
        \end{enumerate}
    \item \textbf{Замена поколения:} Текущая популяция заменяется новой.
    \item \textbf{Проверка условия останова:} Если достигнуто максимальное число поколений, алгоритм завершается, иначе переход к шагу 2.
\end{enumerate}

\begin{figure}[H]
    \centering
    \includegraphics[width=0.7\textwidth]{images/block_diagram_GA.png}
    \caption{Блок-схема генетического алгоритма.}
    \label{fig:ga_flowchart}
\end{figure}

\subsection{Роевой алгоритм}

\textbf{Пошаговое описание:}
\begin{enumerate}
    \item \textbf{Инициализация роя:} создается заданное число частиц, координаты каждой частицы генерируются случайным образом в заданном диапазоне.
    \item \textbf{Инициализация скоростей:} для каждой частицы генерируется значение скорости в заданном диапазоне.
    \item \textbf{Инициализация лучшего локального и глобального значений.}
    \item \textbf{Сохранение значений в историю.} Сохраняются следующие значения: текущее глобальное лучшее значение целевой функции; точка, задающая это лучшее значение; среднее значение целевой функции; положение всех точек (для визуализации).
    \item \textbf{Основной цикл, шаги 5.1 -- 5.4:}
        \begin{enumerate}
            \item \textbf{Генерация коэффициентов направления:} Для каждой частицы и каждой координаты генерируются коэффициенты стремления к локальному и глобальному лучшим значениям.
            \item \textbf{Пересчет скорости} по заданной формуле:
            \[
            v = w v + c_1 r_1 (pbest_x - x) + c_2 r_2 (gbest_x - x)
            \]
            Опционально может быть применено ограничение скорости, в таком случае скорость обрезается в заданном диапазоне.
            \item \textbf{Обновление позиций.}
            \item \textbf{Обновление глобального лучшего значения и обновление истории.}
        \end{enumerate}
    \item \textbf{Условие останова:} алгоритм завершается после выполнения заданного числа итераций. Возвращаются следующие значения: лучшая найденная точка, ее значение; история перемещения лучшей точки; история среднего значения функции по рою; история положений всех частиц (для визуализации).
\end{enumerate}

\begin{figure}[H]
    \centering
    \includegraphics[width=0.7\textwidth]{images/block_diagram_PSO.png}
    \caption{Блок-схема роевого алгоритма.}
    \label{fig:pso_flowchart}
\end{figure}

\section{Формализация задачи}
\label{sec:formalization}

Программа реализована на языке Python с использованием следующих библиотек и фреймворков:
\begin{enumerate}
    \item numpy -- используется для более эффективного выполнения арифметических операций и для генерации случайных величин.
    \item matplotlib -- используется для визуализации работы алгоритмов и для построения графиков.
    \item PyQt6 -- используется для построения графического интерфейса управления параметрами программы.
\end{enumerate}

Исходный код программы имеет следующую многофайловую архитектуру:
\begin{lstlisting}[language=bash, caption={Файловая архитектура программы.}, label=lst:filetree]
|-- ga.py
|-- gui_qt.py
|-- main.py
|-- pso.py
|-- visualisation.py
`-- requirements.txt
\end{lstlisting}

\section{Рекомендации пользователю}
\label{sec:user_guide}

Взаимодействие с программой осуществляется через графический интерфейс.

\section{Рекомендации программиста}
\label{sec:dev_guide}

\textbf{Системные требования:}
\begin{itemize}
    \item Операционная система Windows/Linux/MacOS
    \item Свободное место на диске -- 10 Мб
    \item Свободная оперативная память -- 200 Мб
\end{itemize}

\textbf{Требования к установленному ПО:}
\begin{itemize}
    \item Python 3.9+
    \item numpy, matplotlib, PyQt6. (полный список в файле requirements.txt)
\end{itemize}

\textbf{Инструкция по запуску:}
\begin{enumerate}
    \item Склонировать репозиторий:
    \begin{lstlisting}[language=bash]
git clone https://github.com/pterodactylys/algorithms-and-data-structures.git
    \end{lstlisting}
    \item Перейти в директорию с программой:
    \begin{lstlisting}[language=bash]
cd task_5_ga_and_pso/
    \end{lstlisting}
    \item Активировать виртуальное окружение и установить зависимости:
    \begin{lstlisting}[language=bash]
python3 -m venv new_env
pip install -r requirements.txt
    \end{lstlisting}
    \item Запустить программу:
    \begin{lstlisting}[language=bash]
python3 main.py
    \end{lstlisting}
\end{enumerate}
В результате данных действий откроется Qt-окно с интерфейсом программы.

\section{Описание контрольного примера.}
\label{sec:test_case}

Для проверки корректности работы алгоритма была запущена проверка на задаче оптимизации для функции, описанной выше.

На рисунке ниже изображен интерфейс программы после выполнения алгоритмов. Параметры, с которыми были запущнуты алгоритмы, а также результаты выполнения программы видны на рисунке.

\begin{figure}[H]
    \centering
    \includegraphics[width=0.9\textwidth]{images/interface.png}
    \caption{Интерфейс программы.}
    \label{fig:interface}
\end{figure}

В результате анализа вывода программы приходим к следующим выводам:
\begin{enumerate}
    \item Программа вычисляет значение минимума функции с минимальной погрешностью, разность значений, вычисленных во время работы программы, и значений, известных заранее, является очень маленькой (точность более 12 знаков после запятой). Данная точность достигается как при поиске с помощью ГА, так и с помощью РА. Анализ влияния входных параметров на точность вычисления минимума будет произведен в следующем блоке отчета.
    \item Алгоритмы завершаются за короткий промежуток времени, 0.5 секунды для ГА и 0.005 секунды для РА. Анализ влияния входных параметров на скорость выполнения алгоритмов и сравнение алгоритмов будут произведены в следующем блоке отчета.
\end{enumerate}

Кроме вычисления минимума функции, программа позволяет визуализировать полученные результаты: в программу встроена функциональность по построению графика зависимости значения найденного минимального значения на данной итерации с номером итерации; программа позволяет визуализировать движения особей (частиц) в ходе работы алгоритмов.

\begin{figure}[H]
    \centering
    \includegraphics[width=0.9\textwidth]{images/covergence_plot.png}
    \caption{График зависимости лучшего найденного значения от номера итерации.}
    \label{fig:convergence}
\end{figure}

\begin{figure}[H]
    \centering
    \includegraphics[width=0.9\textwidth]{images/animation_frame.png}
    \caption{Кадр из анимации движения частиц.}
    \label{fig:animation}
\end{figure}

В результате анализа полученных визуальных представлений работы алгоритма приходим к следующим выводам:
\begin{enumerate}
    \item Программа позволяет проводить анализ скорости сходимости алгоритмов в зависимости от входных параметров благодаря построению графика изменения лучшего найденного значения.
    \item Программа позволяет проводить более точный анализ поведения алгоритма в зависимости от входных параметров благодаря удобному инструменту, позволяющему следить за движением особей (частиц) в ходе работы алгоритмов.
\end{enumerate}

\section{Анализ результатов работы алгоритма.}
\label{sec:analysis}

\subsection{Цели анализа}
\begin{enumerate}
    \item Найти оптимальные значения входных параметров для обоих алгоритмов. Оптимальными значениями будем считать те, которые обеспечивают поиск точного значения и при этом не приводят к замедлению работы алгоритма.
    \item Провести анализ алгоритмов в стандартном виде с модифицированными версиями и выявить влияние модификации на характеристики работы алгоритма.
    \item Выявить качественные и количественные изменения в работе алгоритмов при изменении входных значений.
    \item Провести сравнительный анализ работы каждого из алгоритмов, запущенных с оптимальными значениями и определить какой из них является наиболее подходящим для решения задачи оптимизации для имеющейся функции.
\end{enumerate}

\subsection{Методика проведения экспериментов}
Для получения статистически достоверных результатов каждый эксперимент проводился серией из 30 независимых прогонов.
Независимость достигается благодаря изменению параметра генерации случайных величин (seed) при каждом новом запуске.
Данная мера позволяет оценивать не только скорость и точность работы алгоритма, но и его стабильность.

Для упрощения проведения экспериментов, связанных с многократным запуском программы, был реализован скрипт, который
позволяет автоматизировать процесс запуска программы с различными входными параметрами и сохранять полученные
результаты в структурированном виде для последущего анализа.

В качестве основных метрик были использованы:
\begin{itemize}
    \item \textbf{Точность (\(\Delta\))} -- абсолютная ошибка между найденным значением минимума функции и фактическим:
    \[
    \Delta = \left| f_{\text{true}} - f_{\text{найденное}} \right|
    \]
    \item \textbf{Точность по аргументам (\(\Delta_{\arg}\))} -- Евклидово расстояние между найденной точкой минимума и истинной:
    \[
    d = \sqrt{\sum_{i=1}^{n} (x_i^{\text{true}} - x_i^{\text{найденное}})^2}
    \]
    \item \textbf{Стабильность (\(\sigma\))} -- стандартное отклонение ошибки по серии запусков.
    \item \textbf{Скорость сходимости (\(t\))} -- число итераций, необходимое для достижения ошибки менее \(10^{-6}\).
    \item \textbf{Время выполнения (\(\tau\))} -- среднее время работы алгоритма, секунды.
\end{itemize}

Исследование проводилось на области \([-10, 10] \times [-10, 10]\). Размерность пространства \(n = 2\) выбрана в целях упрощения визуализации.

\subsection{Исследование генетического алгоритма}

\subsubsection{Влияние размера популяции}
Размер популяции изменялся от 100 до 1900 с шагом 200.

\begin{table}[H]
    \centering
    \caption{Исследование влияния размера популяции.}
    \resizebox{\textwidth}{!}{%
    \begin{tabular}{c|ccccc}
        \toprule
        \(N\) & \(\Delta\) & \(\Delta_{\arg}\) & \(\sigma\) & \(t\) & \(\tau,\ \text{секунд}\) \\
        \midrule
        100  & \(1.413 \times 10^0\) & \(4.241 \times 10^0\) & \(5.651 \times 10^{-1}\) & \(7.78 \times 10^1\) & \(2.077408 \times 10^{-1}\) \\
        300  & \(5.925 \times 10^{-8}\) & \(4.705 \times 10^{-8}\) & \(5.386 \times 10^{-5}\) & \(5.90 \times 10^1\) & \(6.157380 \times 10^{-1}\) \\
        500  & \(5.100 \times 10^{-8}\) & \(5.860 \times 10^{-8}\) & \(4.542 \times 10^{-5}\) & \(4.09 \times 10^1\) & \(1.015055 \times 10^0\) \\
        700  & \(6.076 \times 10^{-8}\) & \(4.5834 \times 10^{-8}\) & \(5.450 \times 10^{-5}\) & \(3.54 \times 10^1\) & \(1.422419 \times 10^0\) \\
        900  & \(4.573 \times 10^{-8}\) & \(6.392 \times 10^{-8}\) & \(3.836 \times 10^{-5}\) & \(3.70 \times 10^1\) & \(1.826041 \times 10^0\) \\
        1100 & \(2.169 \times 10^{-8}\) & \(2.576 \times 10^{-8}\) & \(3.073 \times 10^{-5}\) & \(3.66 \times 10^1\) & \(2.253967 \times 10^0\) \\
        1300 & \(6.467 \times 10^{-8}\) & \(8.1218 \times 10^{-8}\) & \(5.3285 \times 10^{-5}\) & \(3.68 \times 10^1\) & \(2.640822 \times 10^0\) \\
        1500 & \(5.107 \times 10^{-8}\) & \(5.5472 \times 10^{-8}\) & \(4.662 \times 10^{-5}\) & \(3.54 \times 10^1\) & \(3.135932 \times 10^0\) \\
        1700 & \(2.246 \times 10^{-8}\) & \(1.763 \times 10^{-8}\) & \(3.271 \times 10^{-5}\) & \(3.07 \times 10^1\) & \(3.568018 \times 10^0\) \\
        1900 & \(2.908 \times 10^{-8}\) & \(2.790 \times 10^{-8}\) & \(3.6402 \times 10^{-5}\) & \(2.70 \times 10^1\) & \(4.023912 \times 10^0\) \\
        \bottomrule
    \end{tabular}%
    }
    \label{tab:pop_size}
\end{table}

С увеличением размера популяции увеличивается и время выполнения, зависимость практически линейная. С ростом количества особей повышается точность и стабильность, но начиная со значения в 300 особей прирост становится практически неразличим.

\subsubsection{Влияние вероятности кроссовера}
Вероятность кроссовера изменялась от 0.5 до 1 с шагом 0.1.

\begin{table}[H]
    \centering
    \caption{Исследование влияния вероятности кроссовера.}
    \resizebox{\textwidth}{!}{%
    \begin{tabular}{c|ccccc}
        \toprule
        \(p\) & \(\Delta\) & \(\Delta_{\arg}\) & \(\sigma\) & \(t\) & \(\tau,\ \text{секунд}\) \\
        \midrule
        0.500 & \(2.948971 \times 10^{-8}\) & \(3.772672 \times 10^{-8}\) & \(3.396331 \times 10^{-5}\) & \(3.330000 \times 10^1\) & \(1.954774 \times 10^0\) \\
        0.600 & \(3.372306 \times 10^{-8}\) & \(2.822315 \times 10^{-8}\) & \(4.017700 \times 10^{-5}\) & \(3.380000 \times 10^1\) & \(2.002842 \times 10^0\) \\
        0.700 & \(3.872338 \times 10^{-8}\) & \(4.584356 \times 10^{-8}\) & \(4.059837 \times 10^{-5}\) & \(3.290000 \times 10^1\) & \(2.104592 \times 10^0\) \\
        0.800 & \(1.324232 \times 10^{-8}\) & \(7.873226 \times 10^{-9}\) & \(2.639759 \times 10^{-5}\) & \(3.890000 \times 10^1\) & \(2.093505 \times 10^0\) \\
        0.900 & \(2.826980 \times 10^{-8}\) & \(3.049576 \times 10^{-8}\) & \(3.319865 \times 10^{-5}\) & \(3.510000 \times 10^1\) & \(2.082939 \times 10^0\) \\
        1.000 & \(3.706049 \times 10^{-8}\) & \(2.524861 \times 10^{-8}\) & \(4.373014 \times 10^{-5}\) & \(3.600000 \times 10^1\) & \(2.132884 \times 10^0\) \\
        \bottomrule
    \end{tabular}%
    }
    \label{tab:crossover_prob}
\end{table}

Наилучшие значения по точности достигаются при \(p = 0.8\). В остальных случаях метрики оказались немного хуже.
Вероятно, это связано с тем, что при слишком высокой вероятности кроссовера теряется разнообразие в популяции,
что может привести к преждевременной сходимости к локальному оптимуму.
С другой стороны, при слишком низкой вероятности кроссовера алгоритм может не сохранять достаточное количество
информции о лучших решениях, что также может замедлить процесс поиска оптимума.

\subsubsection{Влияние вероятности мутации}
Вероятность мутации изменялась от 0.05 до 0.5 с шагом 0.05.

\begin{table}[H]
    \centering
    \caption{Исследование влияния вероятности мутации.}
    \resizebox{\textwidth}{!}{%
    \begin{tabular}{c|ccccc}
        \toprule
        \(p\) & \(\Delta\) & \(\Delta_{\arg}\) & \(\sigma\) & \(t\) & \(\tau,\ \text{секунд}\) \\
        \midrule
        0.050 & \(4.894176 \times 10^{-10}\) & \(5.729121 \times 10^{-10}\) & \(4.352485 \times 10^{-6}\) & \(3.950000 \times 10^1\) & \(2.060358 \times 10^0\) \\
        0.100 & \(1.335820 \times 10^{-13}\) & \(1.445321 \times 10^{-13}\) & \(6.330867 \times 10^{-8}\) & \(3.070000 \times 10^1\) & \(2.049491 \times 10^0\) \\
        0.150 & \(1.639933 \times 10^{-11}\) & \(1.729413 \times 10^{-11}\) & \(8.405745 \times 10^{-7}\) & \(2.980000 \times 10^1\) & \(2.113701 \times 10^0\) \\
        0.200 & \(3.052935 \times 10^{-9}\) & \(4.247264 \times 10^{-9}\) & \(1.080024 \times 10^{-5}\) & \(2.960000 \times 10^1\) & \(2.062488 \times 10^0\) \\
        0.250 & \(1.324232 \times 10^{-8}\) & \(7.873226 \times 10^{-9}\) & \(2.639759 \times 10^{-5}\) & \(3.890000 \times 10^1\) & \(2.115668 \times 10^0\) \\
        0.300 & \(1.480270 \times 10^{-7}\) & \(1.472215 \times 10^{-7}\) & \(7.899599 \times 10^{-5}\) & \(6.320000 \times 10^1\) & \(2.093434 \times 10^0\) \\
        0.350 & \(5.048306 \times 10^{-7}\) & \(3.807409 \times 10^{-7}\) & \(1.566273 \times 10^{-4}\) & \(4.830000 \times 10^1\) & \(2.076894 \times 10^0\) \\
        0.400 & \(9.018358 \times 10^{-7}\) & \(8.804333 \times 10^{-7}\) & \(2.059078 \times 10^{-4}\) & \(1.110000 \times 10^2\) & \(2.047508 \times 10^0\) \\
        0.450 & \(1.417778 \times 10^{-6}\) & \(7.413373 \times 10^{-7}\) & \(2.771483 \times 10^{-4}\) & \(7.800000 \times 10^1\) & \(2.056849 \times 10^0\) \\
        0.500 & \(2.012143 \times 10^{-6}\) & \(2.056267 \times 10^{-6}\) & \(2.835861 \times 10^{-4}\) & \(7.600000 \times 10^1\) & \(2.047410 \times 10^0\) \\
        \bottomrule
    \end{tabular}%
    }
    \label{tab:mutation_prob}
\end{table}

Оптимальным оказалось значение \(p = 0.1\). При избыточном увеличении вероятности мутации значительно ухудшается точность поиска оптимума. Помимо этого, растет и количество итераций, необходимых для приближения к окрестности решения.

\subsubsection{Влияние силы мутации}
Коэффициент силы мутации изменялся от 0.1 до 1 с шагом 0.1.

\begin{table}[H]
    \centering
    \caption{Исследование влияния силы мутации.}
    \resizebox{\textwidth}{!}{%
    \begin{tabular}{c|ccccc}
        \toprule
        \(\sigma'\) & \(\Delta\) & \(\Delta_{\arg}\) & \(\sigma\) & \(t\) & \(\tau,\ \text{секунд}\) \\
        \midrule
        0.100 & \(1.034680 \times 10^{-8}\) & \(1.153397 \times 10^{-8}\) & \(2.071156 \times 10^{-5}\) & \(3.150000 \times 10^1\) & \(2.049712 \times 10^0\) \\
        0.200 & \(1.912376 \times 10^{-8}\) & \(1.886803 \times 10^{-8}\) & \(2.898458 \times 10^{-5}\) & \(3.150000 \times 10^1\) & \(2.044004 \times 10^0\) \\
        0.300 & \(7.307934 \times 10^{-8}\) & \(1.226486 \times 10^{-7}\) & \(5.040931 \times 10^{-5}\) & \(3.990000 \times 10^1\) & \(2.041652 \times 10^0\) \\
        0.400 & \(1.352250 \times 10^{-7}\) & \(1.370886 \times 10^{-7}\) & \(7.370266 \times 10^{-5}\) & \(4.350000 \times 10^1\) & \(2.060574 \times 10^0\) \\
        0.500 & \(1.009254 \times 10^{-7}\) & \(1.363793 \times 10^{-7}\) & \(6.544238 \times 10^{-5}\) & \(3.850000 \times 10^1\) & \(2.040837 \times 10^0\) \\
        0.600 & \(1.198760 \times 10^{-7}\) & \(9.666718 \times 10^{-8}\) & \(7.527097 \times 10^{-5}\) & \(5.800000 \times 10^1\) & \(2.039133 \times 10^0\) \\
        0.700 & \(2.741389 \times 10^{-7}\) & \(1.775513 \times 10^{-7}\) & \(1.154844 \times 10^{-4}\) & \(7.480000 \times 10^1\) & \(2.049474 \times 10^0\) \\
        0.800 & \(5.358516 \times 10^{-7}\) & \(6.389252 \times 10^{-7}\) & \(1.475894 \times 10^{-4}\) & \(5.750000 \times 10^1\) & \(2.044652 \times 10^0\) \\
        0.900 & \(5.216789 \times 10^{-7}\) & \(4.946148 \times 10^{-7}\) & \(1.504138 \times 10^{-4}\) & \(8.330000 \times 10^1\) & \(2.095889 \times 10^0\) \\
        1.000 & \(5.309839 \times 10^{-7}\) & \(4.108561 \times 10^{-7}\) & \(1.562627 \times 10^{-4}\) & \(8.430000 \times 10^1\) & \(2.098061 \times 10^0\) \\
        \bottomrule
    \end{tabular}%
    }
    \label{tab:mutation_strength}
\end{table}

При чрезмерном увеличении данного коэффициента увеличивается количество итераций, необходимых для входа в окрестность глобального минимума. Наиболее оптимальным значением оказалось \(\sigma' = 0.1\).

\subsubsection{Влияние турнирного отбора}
Тесты осуществлялись для значений \(k = \{2, 3, 5, 7\}\).

\begin{table}[H]
    \centering
    \caption{Исследование влияния турнирного отбора.}
    \resizebox{\textwidth}{!}{%
    \begin{tabular}{c|ccccc}
        \toprule
        \(k\) & \(\Delta\) & \(\Delta_{\arg}\) & \(\sigma\) & \(t\) & \(\tau,\ \text{секунд}\) \\
        \midrule
        2 & \(1.324232 \times 10^{-8}\) & \(7.873226 \times 10^{-9}\) & \(2.639759 \times 10^{-5}\) & \(3.890000 \times 10^1\) & \(2.074541 \times 10^0\) \\
        3 & \(3.552714 \times 10^{-14}\) & \(7.105427 \times 10^{-15}\) & \(4.218785 \times 10^{-8}\) & \(1.650000 \times 10^1\) & \(2.082557 \times 10^0\) \\
        5 & \(5.684342 \times 10^{-14}\) & \(0\) & \(3.951726 \times 10^{-8}\) & \(1.070000 \times 10^1\) & \(2.092837 \times 10^0\) \\
        7 & \(5.684342 \times 10^{-14}\) & \(0\) & \(4.261332 \times 10^{-8}\) & \(9.400000 \times 10^0\) & \(2.061180 \times 10^0\) \\
        \bottomrule
    \end{tabular}%
    }
    \label{tab:tournament}
\end{table}

С ростом количества особей, участвующих в турнирном отборе, значительно растет точность и скорость поиска минимума. В некоторых случаях данная закономерность может не быть верной, так как алгоритм с высоким значением \(k\) может преждевременно сойтись к локальному минимуму.

\subsubsection{Влияние модификации элитизм}
Параметр принимает два значения: Да / Нет.

\begin{table}[H]
    \centering
    \caption{Исследование влияния элитизма.}
    \resizebox{\textwidth}{!}{%
    \begin{tabular}{c|ccccc}
        \toprule
        Элитизм & \(\Delta\) & \(\Delta_{\arg}\) & \(\sigma\) & \(t\) & \(\tau,\ \text{секунд}\) \\
        \midrule
        Нет & \(9.082491 \times 10^{-6}\) & \(1.003952 \times 10^{-5}\) & \(5.920785 \times 10^{-4}\) & \(5.090000 \times 10^1\) & \(2.075020 \times 10^0\) \\
        Да  & \(1.324232 \times 10^{-8}\) & \(7.873226 \times 10^{-9}\) & \(2.639759 \times 10^{-5}\) & \(3.890000 \times 10^1\) & \(2.083524 \times 10^0\) \\
        \bottomrule
    \end{tabular}%
    }
    \label{tab:elitism}
\end{table}

Полученные выводы имеют схожесть с выводами, сделанными в п.8.3.5. В некоторых случаях модификация элитизм может кратно повысить точность поиска и ускорить его, но также могут быть и случаи, когда модификация приведет к преждевременной сходимости к локальному минимуму.

\subsection{Исследование роевого алгоритма}

\subsubsection{Влияние размера роя}
Количество частиц в рое изменялось от 50 до 500 с шагом 50.

\begin{table}[H]
    \centering
    \caption{Исследование влияния размера роя.}
    \resizebox{\textwidth}{!}{%
    \begin{tabular}{c|ccccc}
        \toprule
        \(N\) & \(\Delta\) & \(\Delta_{\arg}\) & \(\sigma\) & \(t\) & \(\tau,\ \text{секунд}\) \\
        \midrule
        50  & \(3.836931 \times 10^{-13}\) & \(5.237621 \times 10^{-13}\) & \(1.355096 \times 10^{-7}\) & \(8.780000 \times 10^1\) & \(4.569183 \times 10^{-3}\) \\
        100 & \(2.700062 \times 10^{-14}\) & \(1.180442 \times 10^{-14}\) & \(5.342486 \times 10^{-8}\) & \(7.500000 \times 10^1\) & \(5.205083 \times 10^{-3}\) \\
        150 & \(4.547474 \times 10^{-14}\) & \(6.122224 \times 10^{-14}\) & \(5.423153 \times 10^{-8}\) & \(7.030000 \times 10^1\) & \(5.952908 \times 10^{-3}\) \\
        200 & \(2.842171 \times 10^{-14}\) & \(1.797547 \times 10^{-14}\) & \(4.607717 \times 10^{-8}\) & \(6.460000 \times 10^1\) & \(6.808779 \times 10^{-3}\) \\
        250 & \(3.268497 \times 10^{-14}\) & \(1.428173 \times 10^{-14}\) & \(3.795767 \times 10^{-8}\) & \(6.240000 \times 10^1\) & \(7.133583 \times 10^{-3}\) \\
        300 & \(2.842171 \times 10^{-14}\) & \(8.987734 \times 10^{-15}\) & \(4.810519 \times 10^{-8}\) & \(6.170000 \times 10^1\) & \(8.046446 \times 10^{-3}\) \\
        350 & \(3.836931 \times 10^{-14}\) & \(1.109903 \times 10^{-14}\) & \(4.109045 \times 10^{-8}\) & \(5.890000 \times 10^1\) & \(8.533063 \times 10^{-3}\) \\
        400 & \(3.552714 \times 10^{-14}\) & \(9.532931 \times 10^{-15}\) & \(3.861770 \times 10^{-8}\) & \(6.090000 \times 10^1\) & \(9.173525 \times 10^{-3}\) \\
        450 & \(4.405365 \times 10^{-14}\) & \(1.180442 \times 10^{-14}\) & \(3.741379 \times 10^{-8}\) & \(5.800000 \times 10^1\) & \(1.028253 \times 10^{-2}\) \\
        500 & \(3.979039 \times 10^{-14}\) & \(1.238874 \times 10^{-14}\) & \(4.290907 \times 10^{-8}\) & \(5.810000 \times 10^1\) & \(1.104899 \times 10^{-2}\) \\
        \bottomrule
    \end{tabular}%
    }
    \label{tab:swarm_size}
\end{table}

Значительных изменений при изменении количества частиц в рое не происходит. Вероятно, это происходит из-за особенности текущей конфигурации алгоритма, т.е. качество нахождения результата обеспечивается другими параметрами.

\subsubsection{Влияние коэффицента инерции}
Коэффициент инерции изменялся от 0.4 до 1 с шагом 0.1.

\begin{table}[H]
    \centering
    \caption{Исследование влияния коэффициента инерции.}
    \resizebox{\textwidth}{!}{%
    \begin{tabular}{c|ccccc}
        \toprule
        \(w\) & \(\Delta\) & \(\Delta_{\arg}\) & \(\sigma\) & \(t\) & \(\tau,\ \text{секунд}\) \\
        \midrule
        0.400 & \(5.684342 \times 10^{-14}\) & \(4.452869 \times 10^{-8}\) & \(0\) & \(2.030000 \times 10^{1}\) & \(5.707633 \times 10^{-3}\) \\
        0.500 & \(5.684342 \times 10^{-14}\) & \(3.974059 \times 10^{-8}\) & \(0\) & \(2.540000 \times 10^{1}\) & \(5.647646 \times 10^{-3}\) \\
        0.600 & \(5.684342 \times 10^{-14}\) & \(3.678255 \times 10^{-8}\) & \(0\) & \(3.150000 \times 10^{1}\) & \(5.498088 \times 10^{-3}\) \\
        0.700 & \(5.684342 \times 10^{-14}\) & \(4.030899 \times 10^{-8}\) & \(0\) & \(4.260000 \times 10^{1}\) & \(5.457996 \times 10^{-3}\) \\
        0.800 & \(1.861622 \times 10^{-13}\) & \(7.641853 \times 10^{-8}\) & \(3.891517 \times 10^{-13}\) & \(6.740000 \times 10^{1}\) & \(5.456429 \times 10^{-3}\) \\
        0.900 & \(7.138466 \times 10^{-6}\) & \(5.251731 \times 10^{-4}\) & \(8.606383 \times 10^{-6}\) & \(1.558000 \times 10^{2}\) & \(5.320612 \times 10^{-3}\) \\
        1.000 & \(4.448560 \times 10^{-3}\) & \(1.384260 \times 10^{-2}\) & \(4.879384 \times 10^{-3}\) & \(1.103000 \times 10^{2}\) & \(5.277063 \times 10^{-3}\) \\
        \bottomrule
    \end{tabular}%
    }
    \label{tab:inertia}
\end{table}

При избыточном увеличении данного коэффициента наблюдается резкое ухудшение точности поиска минимума функции. Наилучшие результаты достигаются при $w$ в промежутке $[0.4, 0.6]$.
Вероятно, что высокие значения данного параметра приводят к тому, что частицы слишком сильно
ориентируются на свое предыдущее направление движения, что приводит к снижению чувствительности к информации от других частиц и, как следствие, ухужшению качества поиска оптимума.

\subsubsection{Влияние когнитивного коэффициента}
Когнитивный коэффициент изменялся от 0.5 до 2.5 с шагом 0.25.

\begin{table}[H]
    \centering
    \caption{Влияние когнитивного коэффициента.}
    \resizebox{\textwidth}{!}{%
    \begin{tabular}{c|ccccc}
        \toprule
        \(c_1\) & \(\Delta\) & \(\Delta_{\arg}\) & \(\sigma\) & \(t\) & \(\tau,\ \text{секунд}\) \\
        \midrule
        0.500 & \(3.836931 \times 10^{-14}\) & \(9.099387 \times 10^{-15}\) & \(3.875763 \times 10^{-8}\) & \(6.890000 \times 10^1\) & \(5.513012 \times 10^{-3}\) \\
        0.750 & \(5.115908 \times 10^{-14}\) & \(6.961869 \times 10^{-15}\) & \(4.056561 \times 10^{-8}\) & \(6.540000 \times 10^1\) & \(5.406288 \times 10^{-3}\) \\
        1.000 & \(4.263256 \times 10^{-14}\) & \(1.271057 \times 10^{-14}\) & \(5.018416 \times 10^{-8}\) & \(6.480000 \times 10^1\) & \(5.397892 \times 10^{-3}\) \\
        1.250 & \(3.694822 \times 10^{-14}\) & \(6.961869 \times 10^{-15}\) & \(4.394249 \times 10^{-8}\) & \(7.630000 \times 10^1\) & \(5.402979 \times 10^{-3}\) \\
        1.500 & \(1.861622 \times 10^{-13}\) & \(3.891517 \times 10^{-13}\) & \(7.641853 \times 10^{-8}\) & \(6.740000 \times 10^1\) & \(5.382592 \times 10^{-3}\) \\
        1.750 & \(3.930722 \times 10^{-12}\) & \(8.171741 \times 10^{-12}\) & \(3.318538 \times 10^{-7}\) & \(8.350000 \times 10^1\) & \(5.440967 \times 10^{-3}\) \\
        2.000 & \(8.675727 \times 10^{-12}\) & \(1.575956 \times 10^{-11}\) & \(4.958520 \times 10^{-7}\) & \(8.730000 \times 10^1\) & \(5.360579 \times 10^{-3}\) \\
        2.250 & \(2.485440 \times 10^{-9}\) & \(2.507115 \times 10^{-9}\) & \(1.004312 \times 10^{-5}\) & \(1.019000 \times 10^2\) & \(5.339608 \times 10^{-3}\) \\
        2.500 & \(9.477593 \times 10^{-8}\) & \(1.145552 \times 10^{-7}\) & \(5.945898 \times 10^{-5}\) & \(1.363000 \times 10^2\) & \(5.321762 \times 10^{-3}\) \\
        \bottomrule
    \end{tabular}%
    }
    \label{tab:cognitive}
\end{table}

Заметно ухудшение точности при увеличении данного коэффициента. Наилучшие результаты достигаются при
\(c_1\) в промежутке \([0.5, 1.25]\). Вероятно, что высокие значения данного параметра приводят к тому,
что частицы слишком сильно ориентируются на свое предыдущее положение, что может привести к сходимости к локальному минимуму.

\subsubsection{Влияние социального коэффициента.}
Социальный коэффициент изменялся от 0.5 до 2.5 с шагом 0.25.

\begin{table}[H]
    \centering
    \caption{Влияние социального коэффициента.}
    \resizebox{\textwidth}{!}{%
    \begin{tabular}{c|ccccc}
        \toprule
        \(c_2\) & \(\Delta\) & \(\Delta_{\arg}\) & \(\sigma\) & \(t\) & \(\tau,\ \text{секунд}\) \\
        \midrule
        0.500 & \(5.542233 \times 10^{-14}\) & \(4.263256 \times 10^{-15}\) & \(3.955374 \times 10^{-8}\) & \(5.210000 \times 10^1\) & \(5.522650 \times 10^{-3}\) \\
        0.750 & \(5.115908 \times 10^{-14}\) & \(6.961869 \times 10^{-15}\) & \(3.834909 \times 10^{-8}\) & \(5.640000 \times 10^1\) & \(5.595587 \times 10^{-3}\) \\
        1.000 & \(4.831691 \times 10^{-14}\) & \(6.961869 \times 10^{-15}\) & \(3.755681 \times 10^{-8}\) & \(6.380000 \times 10^1\) & \(5.442404 \times 10^{-3}\) \\
        1.250 & \(4.121148 \times 10^{-14}\) & \(7.652779 \times 10^{-15}\) & \(3.969039 \times 10^{-8}\) & \(6.850000 \times 10^1\) & \(5.635683 \times 10^{-3}\) \\
        1.500 & \(1.861622 \times 10^{-13}\) & \(3.891517 \times 10^{-13}\) & \(7.641853 \times 10^{-8}\) & \(6.740000 \times 10^1\) & \(5.524584 \times 10^{-3}\) \\
        1.750 & \(1.820410 \times 10^{-12}\) & \(2.989097 \times 10^{-12}\) & \(2.219393 \times 10^{-7}\) & \(7.900000 \times 10^1\) & \(5.386358 \times 10^{-3}\) \\
        2.000 & \(4.750262 \times 10^{-11}\) & \(5.546970 \times 10^{-11}\) & \(1.325882 \times 10^{-6}\) & \(9.260000 \times 10^1\) & \(5.374771 \times 10^{-3}\) \\
        2.250 & \(3.235499 \times 10^{-10}\) & \(4.022909 \times 10^{-10}\) & \(3.425286 \times 10^{-6}\) & \(1.118000 \times 10^2\) & \(5.358404 \times 10^{-3}\) \\
        2.500 & \(1.663276 \times 10^{-7}\) & \(3.644920 \times 10^{-7}\) & \(6.474510 \times 10^{-5}\) & \(1.154000 \times 10^2\) & \(5.333967 \times 10^{-3}\) \\
        \bottomrule
    \end{tabular}%
    }
    \label{tab:social}
\end{table}

Наилучшие результаты достигаются при \(c_2\) в промежутке \([0.5, 1.25]\).
Вероятно, что высокие значения данного параметра приводят к тому, что частицы
слишком сильно ориентируются на глобальный оптимум, что может привести к сходимости
к локальному оптимуму илик тому, что частицы будут искать глобальный оптимум дольше.

\subsubsection{Влияние коэффициента ограничения скорости.}
Коэффициент ограничения скорости изменялся от 0.05 до 0.5 с шагом 0.05.

\begin{table}[H]
    \centering
    \caption{Влияние коэффициента ограничения скорости.}
    \resizebox{\textwidth}{!}{%
    \begin{tabular}{c|ccccc}
        \toprule
        \(v_{\max}\) & \(\Delta\) & \(\Delta_{\arg}\) & \(\sigma\) & \(t\) & \(\tau,\ \text{секунд}\) \\
        \midrule
        0.050 & \(3.126388 \times 10^{-14}\) & \(1.774935 \times 10^{-14}\) & \(4.453654 \times 10^{-8}\) & \(6.510000 \times 10^1\) & \(5.492954 \times 10^{-3}\) \\
        0.100 & \(2.984279 \times 10^{-14}\) & \(9.947598 \times 10^{-15}\) & \(4.441813 \times 10^{-8}\) & \(6.690000 \times 10^1\) & \(5.402996 \times 10^{-3}\) \\
        0.150 & \(3.552714 \times 10^{-14}\) & \(1.711214 \times 10^{-14}\) & \(5.078329 \times 10^{-8}\) & \(6.050000 \times 10^1\) & \(5.404604 \times 10^{-3}\) \\
        0.200 & \(1.861622 \times 10^{-13}\) & \(3.891517 \times 10^{-13}\) & \(7.641853 \times 10^{-8}\) & \(6.740000 \times 10^1\) & \(5.411475 \times 10^{-3}\) \\
        0.250 & \(2.557954 \times 10^{-14}\) & \(1.530556 \times 10^{-14}\) & \(4.315051 \times 10^{-8}\) & \(7.390000 \times 10^1\) & \(5.390892 \times 10^{-3}\) \\
        0.300 & \(2.984279 \times 10^{-14}\) & \(2.242435 \times 10^{-14}\) & \(5.226198 \times 10^{-8}\) & \(7.610000 \times 10^1\) & \(5.390400 \times 10^{-3}\) \\
        0.350 & \(1.989520 \times 10^{-14}\) & \(1.582454 \times 10^{-14}\) & \(4.567457 \times 10^{-8}\) & \(7.820000 \times 10^1\) & \(5.388371 \times 10^{-3}\) \\
        0.400 & \(3.268497 \times 10^{-14}\) & \(3.868382 \times 10^{-14}\) & \(5.202514 \times 10^{-8}\) & \(8.050000 \times 10^1\) & \(5.381942 \times 10^{-3}\) \\
        0.450 & \(4.405365 \times 10^{-14}\) & \(4.981913 \times 10^{-14}\) & \(6.636644 \times 10^{-8}\) & \(7.710000 \times 10^1\) & \(5.383692 \times 10^{-3}\) \\
        0.500 & \(7.247536 \times 10^{-14}\) & \(1.064297 \times 10^{-13}\) & \(6.914543 \times 10^{-8}\) & \(7.930000 \times 10^1\) & \(5.385762 \times 10^{-3}\) \\
        \bottomrule
    \end{tabular}%
    }
    \label{tab:vmax}
\end{table}

Наилучшие результаты достигаются при \(v_{\max}\) в промежутке \([0.2, 0.3]\). При значениях вне этого промежутка наблюдается
незначительное ухудшение точности поиска оптимума. Вероятно, что слишком высокие значения
данного параметра приводят к слишком быстрому движению частиц, что может привести к тому,
что частицы будут перескакивать через оптимум. С другой стороны, слишком низкие значения данного
параметра могут сказаться в том, что частицы будут двигаться слишком меделнно, что может привести к тому,
что частицы будут дольше искать оптимум, либо к преждевременной сходимости к локальному оптимуму.

Данный результат может быть применен и на случай полного отключения модификации ограничения скорости.
При отключении ограничения по скорости полностью, эффект по "перескакиванию" через оптимум будет еще более заметным.

\subsection{Сравнение генетического и роевого алгоритма}
\subsubsection{Методология сравнения}
Качество оптимизации сравнивалось по следующим метрикам:
\begin{itemize}
    \item \(\Delta\) - абсолютная погрешность найденного решения.
    \item \(\Delta_{\arg}\) - абсолютная погрешность аргумента найденного решения.
    \item \(\sigma\) - стандартное отклонение найденного решения.
    \item \(t\) - количество итераций, необходимых для приближения к окрестности решения.
    \item \(\tau\) - время, затраченное на оптимизацию.
\end{itemize}

Наиболее важными метриками являются \(\Delta\) и \(\Delta_{\arg}\),
так как они напрямую характеризуют качество найденного решения.
Метрика \(\sigma\) позволяет оценить стабильность алгоритма,
а метрики \(t\) и \(\tau\) позволяют оценить эффективность алгоритма.

Выводы о качестве оптимизации будут производиться на основе сравнения этих метрик для обоих алгоритмов, 
запущенных с оптимальными параметрами, определенными в предыдущих пунктах.

\subsubsection{Выбор оптимальных параметров для сравнения}
Для сравнения алгоритмов были выбраны следующие параметры:
\begin{itemize}
    \item Генетический алгоритм (вещественное кодирование): \(N = 200\), \(k = 3\), \(p_{cross} = 0.8\), \(p_{mut} = 0.01\),  \(\sigma' = 0.25\),  элитизм включен, \(e = 3\)
    \item Генетический алгоритм (бинарное кодирование): \(N = 200\), \(k = 3\), \(p_{cross} = 0.8\), \(p_{mut} = 0.25\),  \(\sigma' = 0.25\),  элитизм включен, \(e = 3\)
    \item Роевой алгоритм: \(N = 120\), \(w = 0.4\), \(c_1 = 1.25\), \(c_2 = 1.20\), \(v_{\max} = 0.2\).
\end{itemize}

\subsubsection{Результаты сравнения}
\begin{table}[H]
    \centering
    \caption{Сравнение генетического и роевого алгоритма.}
    \resizebox{\textwidth}{!}{%
    \begin{tabular}{c|ccccc}
        \toprule
        Алгоритм & \(\Delta\) & \(\Delta_{\arg}\) & \(\sigma\) & \(t\) & \(\tau,\ \text{секунд}\) \\
        \midrule
        $GA_{Real}$ & \(5.613288 \times 10^{-14}\) & \(4.066605 \times 10^{-8}\) & \(3.097184 \times 10^{-15}\) & \(1.65 \times 10^1\) & \(5.55 \times 10^{-1}\) \\
        $GA_{Bin}$ & \(0\) & \(2.096648 \times 10^{-9}\) & \(0\) & \(2.35 \times 10^1\) & \(5.44 \times 10^{-1}\) \\        
        $PSO$ & \(5.684342 \times 10^{-14}\) & \(4.270654 \times 10^{-8}\) & \(3.34321 \times 10^{-14}\) & \(1.85 \times 10^1\) & \(1.91 \times 10^{-3}\) \\
        \bottomrule
    \end{tabular}%
    }
    \label{tab:comparison}
\end{table}

Оба алгоритма показали высокую точность поиска оптимума. Погрешность \(\Delta\) и \(\Delta_{\arg}\) оба алгоритма достигли порядка \(10^{-14}\) и \(10^{-8}\) соответственно,
а бинарная версия ГА достигла точности, превышающей точность представления вещественных чисел, что свидетельствует об эффективности обоих алгоритмов в решении данной задачи оптимизации.

Несмотря на схожую точность, роевой алгоритм оказался значительно быстрее генетического алгоритма, что указывает
на его более высокую эффективность в данной задаче. Вероятно, что подход, используемый в роевом алгоритме является
более удачным в данном случае.

Однако, в ходе экспериментов было замечено, что генетический алгоритм оказался более чувсвительным к небольшим
изменениям параметров, что в некоторых случаях приводило к значительному ухудшению качества оптимизации.
С другой стороны, роевой алгоритм показал намного более стабильные результаты  при изменении параметров, что
указывает на его более высокую устойчивость к выбору параметров и стабильность в поиске решения.

\section{Выводы}
В ходе данной лабораторной работы были исследованы и протестированы два алгоритма в задаче оптимизации:
генетический алгоритм и роевой алгоритм. Были проведены эксперименты по исследованию влияния
различных параметров на качество оптимизации, а также проведено сравнение эффективности этих алгоритмов.
Были получены следующие выводы:
\begin{itemize}
    \item Оба алгоритма показали высокую точность в поиске оптимума, что свидетельствует об их применимости для решения задач оптимизации.
    \item Бинарная версия генетического алгоритма ($GA_{Bin}$) показала наилучшую точность среди рассмотренных методов: в серии сравнительных запусков получено \(\Delta = 0\), \(\Delta_{\arg} = 2.096648 \times 10^{-9}\), \(\sigma = 0\).
    \item Для бинарной версии ГА время выполнения составило \(\tau = 5.44 \times 10^{-1}\) с при среднем числе итераций \(t = 2.35 \times 10^1\), что указывает на хорошую вычислительную эффективность при высокой точности.
    \item Роевой алгоритм показал значительно более высокую эффективность по сравнению с генетическим алгоритмом при прочих равных.
    \item Генетический алгоритм оказался менее стабильным в поиске решения и более чувствительным к выбору параметров, что может затрудить подбор этих параметров.
    \item Роевой алгоритм показал более стабильные результаты при изменении параметров.
    \item Выбор оптимальных параметров - наиболее важный этап при решении задач с использованием этих алгоритмов.
    \item Наболее предпочтительным алгоритмом для данной задачи является роевой алгоритм, так как он показал более высокую эффективность и стабильность в поиске решения.
\end{itemize}

Полученные результаты могут распространяться в полном объеме только на оптимизацию
выбранной функции, так как оптимальные параметры алгоритмов могут варьироваться в зависимости
от конкретной задачи. Также, стоит отметить, что результаты могут меняться в зависимости
от начальной инициализации алгоритмов, так как оба алгоритма являются стохастическими.
Тем не менее, проведенное сравнение может показать общие тенденции в поведении алгоритмов
и их чувствительность к различным параметрам.

\pagebreak

\section{Листинг кода}
В данном разделе приведены исходные коды основных модулей программы.

\subsection*{Листинг 2. Исходный код файла \texttt{main.py}}
\VerbatimInput[fontsize=\small, frame=single]{main.py}

\subsection*{Листинг 3. Исходный код файла \texttt{ga.py}}
\VerbatimInput[fontsize=\small, frame=single]{ga.py}

\subsection*{Листинг 4. Исходный код файла \texttt{pso.py}}
\VerbatimInput[fontsize=\small, frame=single]{pso.py}

\subsection*{Листинг 5. Исходный код файла \texttt{gui\_qt.py}}
\VerbatimInput[fontsize=\small, frame=single]{gui_qt.py}

\subsection*{Листинг 6. Исходный код файла \texttt{visualisation.py}}
\VerbatimInput[fontsize=\small, frame=single]{visualisation.py}


\end{document}
